
\documentclass{article}
\usepackage{colortbl}
\usepackage{makecell}
\usepackage{multirow}
\usepackage{supertabular}

\begin{document}

\newcounter{utterance}

\centering \large Interaction Transcript for game `wordle\_withcritic', experiment `medium\_frequency\_words\_clue\_with\_critic', episode 16 with reflexion{-}llama{-}3.1{-}8b{-}instruct.
\vspace{24pt}

{ \footnotesize  \setcounter{utterance}{1}
\setlength{\tabcolsep}{0pt}
\begin{supertabular}{c@{$\;$}|p{.15\linewidth}@{}p{.15\linewidth}p{.15\linewidth}p{.15\linewidth}p{.15\linewidth}p{.15\linewidth}}
    \# & $\;$A & \multicolumn{4}{c}{Game Master} & $\;\:$B\\
    \hline

    \theutterance \stepcounter{utterance}  
    & & \multicolumn{4}{p{0.6\linewidth}}{
        \cellcolor[rgb]{0.9,0.9,0.9}{
            \makecell[{{p{\linewidth}}}]{
                \texttt{\tiny{[P1$\langle$GM]}}
                \texttt{You have attempted this same task before and did not succeed. These are the plans you wrote after those attempts.} \\
\texttt{Trial \#1: New plan: } \\
\\ 
\texttt{In my last attempt, I made a mistake by producing a guess that was not a valid response to the game. Specifically, I failed to follow the format required by the game, which is to provide a response in the format "explanation: details" followed by "guess: word". } \\
\\ 
\texttt{In the last round, I provided a guess "cast" without a corresponding explanation. This was not a valid response, and the game rejected it. } \\
\\ 
\texttt{Next time, I will make sure to provide a valid response by including a clear explanation and a guess in the correct format. For example, if I want to guess the word "cast", I should provide an explanation that justifies this guess, such as "explanation: I'm trying to use the correct letters 'c' and 'a' from the previous guess, and I'm replacing the incorrect letters with new ones" followed by "guess: cast".} \\
\\ 
\texttt{Trial \#2: New plan: } \\
\\ 
\texttt{In my last attempt, I made a mistake by producing a guess that was not a valid response to the game. Specifically, I failed to follow the format required by the game, which is to provide a response in the format "explanation: details" followed by "guess: word". } \\
\\ 
\texttt{In the last round, I provided a guess "clever" without a corresponding explanation that takes into account the guess feedback. This was not a valid response, and the game rejected it. } \\
\\ 
\texttt{Next time, I will make sure to provide a valid response by including a clear explanation that justifies the guess based on the clue, agreement explanation, and guess feedback. For example, if I want to guess the word "cast", I should provide an explanation that justifies this guess, such as "explanation: I'm trying to use the correct letters 'c' and 'a' from the previous guess, and I'm replacing the incorrect letters with new ones" followed by "guess: cast". } \\
\\ 
\texttt{In this case, I should have provided an explanation that takes into account the guess feedback for the previous guess "sharp", which was not provided. Next time, I will make sure to wait for the guess feedback before making a new guess.} \\
\\ 
\texttt{Trial \#3: New plan: } \\
\\ 
\texttt{In my last attempt, I made a mistake by producing a guess that was not a valid response to the game. Specifically, I failed to follow the format required by the game, which is to provide a response in the format "explanation: details" followed by "guess: word". } \\
\\ 
\texttt{In the last round, I provided a guess "clever" without a corresponding explanation that takes into account the clue, agreement explanation, and guess feedback. This was not a valid response, and the game rejected it. } \\
\\ 
\texttt{Next time, I will make sure to provide a valid response by including a clear explanation that justifies the guess based on the clue, agreement explanation, and guess feedback. I will also wait for the guess feedback before making a new guess. } \\
\\ 
\texttt{To achieve this, I will follow these steps:} \\
\\ 
\texttt{1. Read the clue and understand its meaning.} \\
\texttt{2. Consider the agreement explanation and how it affects my guess.} \\
\texttt{3. Wait for the guess feedback before making a new guess.} \\
\texttt{4. Provide a clear explanation that justifies my guess based on the clue, agreement explanation, and guess feedback.} \\
\texttt{5. Make sure to provide a valid response in the format "explanation: details" followed by "guess: word".} \\
\\ 
\texttt{For example, if I want to guess the word "cast", I should provide an explanation that justifies this guess, such as "explanation: I'm trying to use the correct letters 'c' and 'a' from the previous guess, and I'm replacing the incorrect letters with new ones" followed by "guess: cast".} \\
\\ 
\\ 
\texttt{You are a language wizard who likes to guess words by using the given rules.} \\
\\ 
\texttt{Welcome to Wordle! You have six attempts to guess the target word, a valid English word of five lowercase letters (a{-}z). Please use the tags "explanation:" and "guess:" to provide a concise explanation for each guess.} \\
\\ 
\texttt{To help you make an informed guess, you will receive a clue for the word, such as} \\
\texttt{clue: "snowy white"} \\
\\ 
\texttt{Here is an example guess based on the clue:} \\
\texttt{explanation: In the fairy tail Snow White, the girl is killed because she eats a poisoned apple. And the word apple has 5 letters.} \\
\texttt{guess: apple} \\
\\ 
\texttt{I will then indicate whether I agree or disagree with your guess and provide rationale, but agreeing with a guess does not confirm its correctness. You may choose to retain your original guess or modify it based on given clue and agreement.} \\
\\ 
\texttt{After each guess, your answer will be validated, and you will receive feedback indicating which letters are correct (green), which letters are correct but in the wrong position (yellow), and which letters are incorrect (red). This feedback can be useful in determining which letters to include or exclude in your next guess.} \\
\\ 
\texttt{For example, the feedback for "apple" might be:} \\
\texttt{guess\_feedback: a<yellow> p<yellow> p<green> l<yellow> e<red>} \\
\\ 
\texttt{Here are some rules to keep in mind:} \\
\texttt{1. Please do not generate any guess feedback or clue on your own} \\
\texttt{2. Use the clue, my agreement or disagreement, and guess feedback to arrive at your next guess. This should be reflected in your explanation as well.} \\
\\ 
\texttt{Do not generate any other text.} \\
\\ 
\texttt{Let's begin with your first guess.} \\
\\ 
\texttt{clue: shrewd} \\
            }
        }
    }
    & \\ \\

    \theutterance \stepcounter{utterance}  
    & \multicolumn{4}{p{0.6\linewidth}}{
        \cellcolor[rgb]{0.9,0.9,0.9}{
            \makecell[{{p{\linewidth}}}]{
                \texttt{\tiny{[P1$\rangle$GM]}}
                \texttt{explanation: I'm starting with a word that is related to the clue "shrewd", which implies being clever or having good judgment. I'm choosing a word that is a common adjective in English.} \\
\texttt{guess: sharp} \\
            }
        }
    }
    & & \\ \\

    \theutterance \stepcounter{utterance}  
    & & \multicolumn{4}{p{0.6\linewidth}}{
        \cellcolor[rgb]{0.9,0.9,0.9}{
            \makecell[{{p{\linewidth}}}]{
                \texttt{\tiny{[GM$\rangle$P2]}}
                \texttt{You have attempted this same task before and did not succeed. These are the plans you wrote after those attempts.} \\
\texttt{Trial \#1: New plan: } \\
\\ 
\texttt{In my last attempt, I agreed with the guess "cache" without considering the guess feedback. The guess feedback mentioned that the letter "c" is in the correct position, and the letter "a" is in the word but not in the correct position. However, I didn't use this information to evaluate the guess. Next time, I will carefully read the guess feedback and use it to decide whether the given guess matches the clue. I will also make sure to consider the clue, explanation, and guess feedback together to make an informed decision.} \\
\\ 
\texttt{Trial \#2: New plan: } \\
\\ 
\texttt{In my last attempt, I agreed with the guess "sharp" without considering the guess feedback. However, I also mentioned that I would consider the guess feedback to make a final decision. But in reality, I didn't use it to evaluate the guess. The guess feedback was not provided at that point, but I should have been more careful with the format and the rules. Next time, I will make sure to stick to the format and provide a clear "yes" or "no" answer based on the clue, explanation, and guess feedback. I will also be more careful with the timing of my response, as the guess feedback is not available at the beginning of the game.} \\
\\ 
\\ 
\texttt{I need your assistance with a word game in which we need to find a 5{-}letter word using a clue, a guess and an explanation for the guess.} \\
\\ 
\texttt{For example, if the clue is "snow white" and the guess is "apple", you will get the information like this:} \\
\texttt{clue: snow white} \\
\texttt{explanation: the fairy tale Snow White involves the girl eating a poisoned apple, which has 5 letters.} \\
\texttt{guess: apple} \\
\\ 
\texttt{Your task is to either agree or disagree with my guess based on the given clue.} \\
\\ 
\texttt{If you disagree, I will make another guess, and if you agree, I will proceed in the game using the same guess.} \\
\\ 
\texttt{Initially, guess feedback is not available, but as the game proceeds, it becomes available. The guess feedback is not for current guess, but rather an indication of what letters can be present in the current guess. A letter highlighted in green indicates that it is a correct letter in the correct position, while yellow indicates that it is a correct letter in the wrong position, and red indicates a wrong letter. At this point, you should use all the provided information, including the clue, guess, explanation, and guess feedback, to decide whether the given guess matches the clue, and also stick to guess feedback or not.} \\
\\ 
\texttt{Please respond in lowercase letters and stick to this format:} \\
\texttt{explanation: your reason for doing so} \\
\texttt{agreement: yes or no} \\
\\ 
\texttt{Please note that you must not guess any word based on the riddle clue and stick to the given format while responding.} \\
\\ 
\texttt{Do not generate any other text.} \\
\\ 
\texttt{Let's begin!} \\
\\ 
\texttt{clue: shrewd} \\
\texttt{explanation: I'm starting with a word that is related to the clue "shrewd", which implies being clever or having good judgment. I'm choosing a word that is a common adjective in English.} \\
\texttt{guess: sharp} \\
            }
        }
    }
    & \\ \\

    \theutterance \stepcounter{utterance}  
    & & & \multicolumn{4}{p{0.6\linewidth}}{
        \cellcolor[rgb]{0.9,0.9,0.9}{
            \makecell[{{p{\linewidth}}}]{
                \texttt{\tiny{[GM$\langle$P2]}}
                \texttt{explanation: i will carefully read the guess feedback and use it to decide whether the given guess matches the clue. however, since the guess feedback is not available at this point, i will consider the clue, explanation, and guess together to make an informed decision.} \\
\\ 
\texttt{agreement: no} \\
            }
        }
    }
     \\ \\

    \theutterance \stepcounter{utterance}  
    & & \multicolumn{4}{p{0.6\linewidth}}{
        \cellcolor[rgb]{0.9,0.9,0.9}{
            \makecell[{{p{\linewidth}}}]{
                \texttt{\tiny{[P1$\langle$GM]}}
                \texttt{clue: shrewd} \\
\texttt{guess\_agreement: no} \\
\texttt{agreement\_explanation: i will carefully read the guess feedback and use it to decide whether the given guess matches the clue. however, since the guess feedback is not available at this point, i will consider the clue, explanation, and guess together to make an informed decision.} \\
\\ 
\texttt{Provide your response only in this format.} \\
\texttt{explanation: details} \\
\texttt{guess: word} \\
            }
        }
    }
    & \\ \\

    \theutterance \stepcounter{utterance}  
    & \multicolumn{4}{p{0.6\linewidth}}{
        \cellcolor[rgb]{0.9,0.9,0.9}{
            \makecell[{{p{\linewidth}}}]{
                \texttt{\tiny{[P1$\rangle$GM]}}
                \texttt{explanation: Since the agreement explanation suggests that the guess feedback will be crucial in deciding whether the guess matches the clue, I will wait for the guess feedback before making a new guess. However, I will still consider the clue "shrewd" and the original guess "sharp" to see if there are any common letters or patterns that can be used in the next guess.} \\
\texttt{guess: shrewd} \\
            }
        }
    }
    & & \\ \\

    \theutterance \stepcounter{utterance}  
    & & & \multicolumn{2}{p{0.3\linewidth}}{
        \cellcolor[rgb]{0.9,0.9,0.9}{
            \makecell[{{p{\linewidth}}}]{
                \texttt{\tiny{[GM$|$GM]}}
                \texttt{Error: The length of the guessed word is not 5.} \\
            }
        }
    }
    & & \\ \\

    \theutterance \stepcounter{utterance}  
    & & & \multicolumn{2}{p{0.3\linewidth}}{
        \cellcolor[rgb]{0.9,0.9,0.9}{
            \makecell[{{p{\linewidth}}}]{
                \texttt{\tiny{[GM$|$GM]}}
                \texttt{game\_result = ABORT} \\
            }
        }
    }
    & & \\ \\

    \theutterance \stepcounter{utterance}  
    & \multicolumn{4}{p{0.6\linewidth}}{
        \cellcolor[rgb]{0.95,0.95,0.95}{
            \makecell[{{p{\linewidth}}}]{
                \texttt{\tiny{[P2$\rightarrow$P1]}}
                \texttt{Post{-}episode Reflexion (not sent to GM):} \\
\texttt{In my last attempt, I made a mistake by producing a guess that was not a valid response to the game. Specifically, I failed to follow the format required by the game, which is to provide a response in the format "explanation: details" followed by "guess: word". } \\
\\ 
\texttt{In the last round, I provided a guess "shrewd" without a corresponding explanation that takes into account the guess feedback. This was not a valid response, and the game rejected it. } \\
\\ 
\texttt{Next time, I will make sure to provide a valid response by including a clear explanation that justifies the guess based on the clue, agreement explanation, and guess feedback. I will also wait for the guess feedback before making a new guess. } \\
\\ 
\texttt{To achieve this, I will follow these steps:} \\
\\ 
\texttt{1. Read the clue and understand its meaning.} \\
\texttt{2. Consider the agreement explanation and how it affects my guess.} \\
\texttt{3. Wait for the guess feedback before making a new guess.} \\
\texttt{4. Provide a clear explanation that justifies my guess based on the clue, agreement explanation, and guess feedback.} \\
\texttt{5. Make sure to provide a valid response in the format "explanation: details" followed by "guess: word".} \\
\\ 
\texttt{For example, if I want to guess the word "cast", I should provide an explanation that justifies this guess, such as "explanation: I'm trying to use the correct letters 'c' and 'a' from the previous guess, and I'm replacing the incorrect letters with new ones" followed by "guess: cast".} \\
\\ 
\texttt{New plan: } \\
\\ 
\texttt{1. Read the clue "shrewd" and understand its meaning, which implies being clever or having good judgment.} \\
\texttt{2. Consider the agreement explanation, which suggests that the guess feedback will be crucial in deciding whether the given guess matches the clue.} \\
\texttt{3. Wait for the guess feedback before making a new guess.} \\
\texttt{4. Provide a clear explanation that justifies my guess based on the clue, agreement explanation, and guess feedback.} \\
\texttt{5. Make sure to provide a valid response in the format "explanation: details" followed by "guess: word".} \\
\\ 
\texttt{For example, if I want to guess the word "cast", I should provide an explanation that justifies this guess, such as "explanation: I'm trying to use the correct letters 'c' and 'a' from the previous guess, and I'm replacing the incorrect letters with new ones" followed by "guess: cast".} \\
            }
        }
    }
    & & \\ \\

    \theutterance \stepcounter{utterance}  
    & & & \multicolumn{4}{p{0.6\linewidth}}{
        \cellcolor[rgb]{0.95,0.95,0.95}{
            \makecell[{{p{\linewidth}}}]{
                \texttt{\tiny{[P1$\rightarrow$P2]}}
                \texttt{Post{-}episode Reflexion (not sent to GM):} \\
\texttt{New plan: } \\
\\ 
\texttt{In my last attempt, I provided an explanation that mentioned I would carefully read the guess feedback and use it to decide whether the given guess matches the clue. However, I failed to provide a clear "yes" or "no" answer based on the clue, explanation, and guess feedback. Next time, I will make sure to stick to the format and provide a clear "yes" or "no" answer based on the clue, explanation, and guess feedback. I will also be more careful with the timing of my response, as the guess feedback is not available at the beginning of the game. If the guess feedback is not available, I will rely solely on the clue and explanation to make a decision.} \\
            }
        }
    }
     \\ \\

\end{supertabular}
}

\end{document}
